% Part: many-valued-logic
% Chapter: tableaux 
% Section: common-rules

\documentclass[../../../include/open-logic-section]{subfiles}

\begin{document}

\olfileid{mvl}{tab}{int}

\olsection{Introduction}

\subsection{Setting up the tableaux}

A tableau allows us to work out, from a set of hypotheses, whether 
there are !!{valuation}s that would satisfy these hypotheses. If all
the branche close, we know that there is no !!{valuation}, hence no
!!{structure}, corresponding to these hypotheses.

To show that $\Gamma \Entails[\Log{X}] !A$, we make the hypothesis that
it is \emph{not} the case, and we use a tableau to work out whether some 
!!{valuation} would be a counter-example to the consequence. 

In bivalent logic, the hypothesis to examine is simple: premises true, 
conclusion false. 

In trivalent logic, however, the hypotheses is: premises \emph{of the 
designated value(s)}, conclusion \emph{of the undesignated value}. In
three-valued logics, 
\begin{itemize}
	\item In truth-preservation logics such as $\LogKw$, $\LogKs$, $\LogLuk[3]$ ($\True$ 
	is the sole designated value), we should consider the hypothesis 
	that the premise is $\True$ and the conclusion $\False$ \emph{or $\Undef$}.
	\item In non-falsehood-preservation logics such as $\LogLP$ ($\True, \Undef$ are the 
	designated values), we should consider the hypothesis that the premises 
	are $\True$ \emph{or $\Undef$} and the conclusion false. We need to consider
	\emph{each} possible combination of $\True$ and $\Undef$ for the premises. For
	two premises, that is four possibilities (both $\True$, first $\True$ second $\Undef$,
	first $\Undef$ second $\True$, both $\Undef$); for three, nine, etc. 
\end{itemize}

For example, this is how to begin a tableau for $!A,\lnot !A \lor !B \Entails !B$ 
in $\LogLuk[3]$ or $\LogKs$:

\begin{oltableau}
	[\sFmla{\True}{\formula{A}}, just=\TAss
		[\sFmla{\True}{\lnot \formula{A} \lor \formula{B}}, just=\TAss
			[\sFmla{\False}{\formula{B}}, just=\TAss]
			[\sFmla{\Undef}{\formula{B}}, just=\TAss]
		]
	]
\end{oltableau}

On the other hand, the \emph{development} of the tableau only depends on
the semantics. Therefore:
\begin{enumerate}
	\item The rules for $\LogKs$ and $\LogLP$ are the same. They only differ
	in the hypotheses used to set up a tableau.
	\item $\LogKs$ and $\LogLuk[3]$ share the rules for $\lnot$, $\land$, $\lor$.
	\item All three-valued logics share the rule for $\lnot$.
\end{enumerate}

$\LogLuk[3]$ sets up its tableaux like $\LogKs$, but uses a different rule for the 
conditional. $\LogLP$ uses the very same rules as $\LogKs$, but sets up its 
tableaux differently. The rules are determined by the semantics alone; the set-up
is determined by which values the logic treat as designated.

\subsection{Signed tableaux}

In bivalent logic, we can use either \emph{signed} tableaux, 
in which we consider (for each main operator) the case where !!a{formula}
(with that main operator) is true, or false, or \emph{unsigned} tableaux,
in which we consider instead (for each main operator) !!a{formula}
(with that main operator) itself and its negation. In many-valued logic,
we have three or more cases to consider, e.g. those where !!a{formula}
is $\True$, $\False$, and $\Undef$. This requires us to use \emph{signed}
tableaux.



\end{document}
